% Fragmento integrable. Preambulo anfitrion: amsmath, amssymb, longtable, hyperref.
\section{Clasificación y continuación del selector de primeras raíces}
\label{hmtfg:fragmento:clasificacion-y-continuacion-del-selector-de-primeras-raices}

\subsection{Objeto y procedencia}
\label{hmtfg:clasificacion:1}

Este apéndice conserva la familia publicada por el lector dodecafásico HMT:

\[
u_0(x)=\frac1{12}\sum_{m=1}^{12}\left(1-\cos\frac{2(3m-1)x}{\sqrt{899}}\right),
\qquad f_\lambda(x)=1-\lambda u_0(x).
\]

La familia procede de los antecedentes operatorios y analíticos expuestos en este desarrollo. La existencia y la compacidad de sus realizaciones del itinerario infinito se han demostrado en las subsecciones de realización de la palabra y conservación de sus signos.

El argumento siguiente es propio y utiliza exclusivamente esas propiedades y las operaciones de retorno de la familia fijada. No utiliza como entrada un parámetro de otra familia ni el valor de una constante universal. Las hojas de los itinerarios aquí utilizados son las dos ramas monótonas del lector crítico; las restantes coordenadas y la memoria HMT conservan su residencia en el estado del que procede el lector.

Escribimos
\[
c_j(\lambda)=f_\lambda^j(0),\qquad
\tau_j=(-1)^{v_2(j)},\qquad
W_n=\tau_1\cdots\tau_{2^n-1}.
\]
El itinerario de una órbita crítica que retorna por primera vez al origen termina en el símbolo cero. El orden de itinerarios es el orden unimodal de máximo: antes de comparar el símbolo de índice \(j\), el factor de orientación es
\[
O_{j-1}(K)=\prod_{i=1}^{j-1}(-K_i).
\]
El orden ordinario de los símbolos es \(-1<0<1\).


\subsection{Clasificación anterior a la palabra infinita}
\label{hmtfg:clasificacion:2}

\textbf{Teorema de clasificación.} Sea \(F:I\to I\), con \(I=[\ell,1]\), \(\ell<0\), un mapa continuo, estrictamente creciente a la izquierda de cero y estrictamente decreciente a su derecha, cuyo máximo satisface \(F(0)=1\). Supóngase que cero tiene período exacto \(2^n\), con \(n\ge1\), y que su itinerario crítico \(K\) satisface \(K<\tau\). Entonces
\[
\boxed{K=W_n0.}
\]
La conclusión se refiere al itinerario, y permite que parámetros distintos de una familia tengan ese mismo itinerario.

\textbf{Demostración.} Denotemos \(d_j=F^j(0)\). Para período dos el itinerario es \(+0=W_10\). Consideremos un período exacto \(p=2^n\ge4\).

Se tiene \(d_2<0\). En efecto, \(d_2=0\) daría período dos. Si \(d_2>0\), la monotonía de la rama positiva da
\[
F([d_2,1])=[d_2,d_3]\subset[d_2,1],
\]
y la órbita crítica permanece en un intervalo estrictamente positivo, incompatible con el retorno a cero.

Los signos iniciales son entonces \(+,-\). El factor de orientación para comparar el tercer signo es \(-1\); por ello \(K<\tau\) obliga a \(d_3>0\). Un cero en ese lugar produciría período tres.

Tampoco puede ocurrir \(d_4<0\). Bajo esta hipótesis,
\[
d_2<d_4<0,
\qquad F([d_2,d_4])=[d_3,d_5]\subset(0,1],
\]
pues \(d_4=F(d_3)>F(1)=d_2\) y \(d_5=F(d_4)>F(d_2)=d_3\). La rama positiva es decreciente, de modo que
\[
F^2([d_2,d_4])=[d_6,d_4]\subset[d_2,d_4].
\]
Aquí \(d_6\ge d_2\), porque \(d_5\le1\), y \(d_6<d_4\), porque \(d_5>d_3\). Este intervalo estrictamente negativo atrapa los retornos pares de la órbita crítica, y excluye el retorno a cero. Por tanto \(d_4\ge0\). Si \(p=4\), se obtiene precisamente \(+-+0=W_20\).

Supongamos ahora \(p\ge8\). Entonces \(d_4>0\). Los factores de orientación para los signos quinto y sexto son, respectivamente, \(-1\) y \(+1\). Como \(\tau_5=+1\) y \(\tau_6=-1\), la desigualdad \(K<\tau\) obliga a
\begin{equation}
\operatorname{sign}(d_1,\ldots,d_6)=(+,-,+,+,+,-).
\label{hmtfg:clasificacion:eq:1}
\end{equation}
Los signos cero en esos lugares están excluidos por el período exacto.

Construimos el intervalo de retorno
\[
J=[d_2,d_4].
\]
La desigualdad \(F(d_3)=d_4>0>d_6=F(d_5)\), junto con \(d_3,d_5>0\), implica \(d_3<d_5\); por tanto
\[
F(J)=[d_3,1].
\]
Además \(d_4<d_3\). Para comprobarlo, la imagen por \(F\) del intervalo de extremos \(d_3,d_4\) es estrictamente positiva, pues sus imágenes extremas son \(d_4,d_5>0\). En ese intervalo \(F^2\) es estrictamente creciente. Sus valores extremos satisfacen
\[
F^2(d_3)=d_5>0>d_6=F^2(d_4),
\]
lo que fuerza \(d_3>d_4\). En consecuencia,
\begin{equation}
J\cap F(J)=\varnothing,
\qquad F^2(J)=F([d_3,1])=[d_2,d_4]=J.
\label{hmtfg:clasificacion:eq:2}
\end{equation}

El mapa \(F^2|_J\) posee un único mínimo en cero. La conjugación por \(h(x)=d_2x\), que invierte la orientación, define
\begin{equation}
G=h^{-1}\circ F^2\circ h:
\left[\frac{d_4}{d_2},1\right]\longrightarrow
\left[\frac{d_4}{d_2},1\right].
\label{hmtfg:clasificacion:eq:3}
\end{equation}
Este mapa satisface las mismas hipótesis de máximo que \(F\), con \(G(0)=1\), y su crítico tiene período exacto \(p/2\). Su itinerario es
\begin{equation}
K'_j=-K_{2j}.
\label{hmtfg:clasificacion:eq:4}
\end{equation}
Todos los símbolos impares de \(K\) son positivos, porque las posiciones pares pertenecen a \(J\) y las impares a \(F(J)\subset(0,1]\).

La transformación \eqref{hmtfg:clasificacion:eq:4} conserva el orden de comparación con \(\tau\). En efecto, \(\tau_{2j-1}=+1\) y \(\tau_{2j}=-\tau_j\). Si \(q\) es el primer lugar en que \(K'\) y \(\tau\) difieren, el primer lugar de diferencia entre \(K\) y \(\tau\) es \(2q\), y
\begin{equation}
O_{2q-1}(K)=-O_{q-1}(K'),
\qquad K_{2q}-\tau_{2q}=-(K'_q-\tau_q).
\label{hmtfg:clasificacion:eq:5}
\end{equation}
El producto que determina el orden es idéntico. Las identidades siguen siendo válidas cuando el primer símbolo distinto es el cero final. Así \(K'<\tau\).

La inducción aplicada a \(G\) da \(K'=W_{n-1}0\). Las posiciones impares positivas y \eqref{hmtfg:clasificacion:eq:4} reconstruyen exactamente \(K=W_n0\). Queda demostrado el teorema. \(\square\)

El teorema demuestra también que cada mapa clasificado posee los retornos restrictivos de duplicación anteriores a su retorno superestable, con los dominios de \eqref{hmtfg:clasificacion:eq:2}–\eqref{hmtfg:clasificacion:eq:3}. No presupone una monotonía del parámetro de una familia.

\textbf{Normalización simétrica de la familia fijada.} Cuando \(F\) es par y está definido sobre \([-1,1]\), las desigualdades anteriores dan además
\[
F(d_4)=d_5>d_3=F(d_2)=F(-d_2).
\]
Como ambos argumentos \(d_4,-d_2\) son positivos y \(F\) decrece en esa rama,
\begin{equation}
0<d_4<-d_2<1.
\label{hmtfg:clasificacion:eq:5a}
\end{equation}
Por tanto \(F^2([d_2,-d_2])=[d_2,d_4]\), y la misma conjugación \(h(x)=d_2x\) define sobre todo \([-1,1]\) un mapa par de máximo cuya imagen es \([d_4/d_2,1]\subset[-1,1]\). Coincide exactamente con el operador \(RF=-F(F(-ax))/a\), \(a=-F(1)\), utilizado en el corpus. Esta ampliación simétrica conserva el retorno crítico, su período y su itinerario.


\subsection{Primer parámetro de itinerario infinito}
\label{hmtfg:clasificacion:3}

Trabajamos sobre el intervalo compacto ya utilizado en la demostración de existencia,
\[
L=\left[\frac1{2u_0(1)},\frac2{u_0(1)}\right].
\]
Sea
\[
\Lambda_\tau=\{\lambda\in L:
\operatorname{sign}c_j(\lambda)=\tau_j\text{ para todo }j\ge1\}.
\]
Por las demostraciones de existencia y separación de signos de la prolongación, este conjunto es no vacío y compacto. Por tanto existe
\begin{equation}
a=\min\Lambda_\tau.
\label{hmtfg:clasificacion:eq:6}
\end{equation}

\textbf{Lema de comparación.} Para todo \(\lambda\in[1/(2u_0(1)),a)\), se tiene \(K_\lambda<\tau\).

\textbf{Demostración.} En el extremo izquierdo, el itinerario es \(+^\infty<\tau\). El argumento de separación de la realización de la palabra infinita demuestra que los conjuntos de comparación estricta con \(\tau\) son abiertos: en un itinerario infinito distinto de \(\tau\), la primera diferencia finita persiste por continuidad; en un parámetro superestable, las dos palabras periódicas vecinas y la palabra con cero final se sitúan en el mismo lado de \(\tau\), por su maximalidad estricta. El intervalo anterior a \(a\) carece, por definición, de parámetros con itinerario \(\tau\). Su conexidad mantiene la comparación inicial en todo ese intervalo. \(\square\)


\subsection{Existencia de cada primera raíz nueva}
\label{hmtfg:clasificacion:4}

Conservamos el selector original. La primera raíz es
\[
\lambda_1=\frac1{u_0(1)},
\]
y, para \(n\ge2\), \(\lambda_n\) es la primera raíz de \(c_{2^n}\) situada a la derecha de \(\lambda_{n-1}\) que satisface
\begin{equation}
c_{2^j}(\lambda_n)\ne0\quad(1\le j<n).
\label{hmtfg:clasificacion:eq:7}
\end{equation}
Como cualquier primer período que divida \(2^n\) es una potencia de dos, \eqref{hmtfg:clasificacion:eq:7} equivale a exigir período crítico exacto \(2^n\).

\textbf{Lema de existencia.} Todas las raíces de este selector existen y cumplen
\begin{equation}
\lambda_1<\lambda_2<\cdots<\lambda_n<a.
\label{hmtfg:clasificacion:eq:8}
\end{equation}

\textbf{Demostración.} El signo \(c_2(a)<0\) implica \(\lambda_1<a\). Supongamos construida \(\lambda_{n-1}<a\), y pongamos \(p=2^{n-1}\). La función analítica \(c_p(\lambda)\) no es idénticamente cero, puesto que \(c_p(a)\ne0\). Su conjunto de ceros en \([\lambda_{n-1},a]\) es finito, no vacío y carece de \(a\). Sea \(z<a\) su último cero.

En \((z,a]\), el signo de \(c_p\) es constante e igual a \(s=\tau_p\). Para \(g_\lambda=f_\lambda^p\), se tiene
\[
g_\lambda(0)=c_p(\lambda),\qquad g'_\lambda(0)=0.
\]
Una cota uniforme local de la segunda derivada da
\begin{equation}
c_{2p}(\lambda)
=g_\lambda(c_p(\lambda))
=c_p(\lambda)+O(c_p(\lambda)^2).
\label{hmtfg:clasificacion:eq:9}
\end{equation}
Al aproximarse \(\lambda\) a \(z\) por la derecha, el término de primer orden domina; por tanto \(c_{2p}(\lambda)\) tiene signo \(s\). En \(a\), su signo es
\[
\tau_{2p}=-\tau_p=-s.
\]
El teorema del valor intermedio proporciona un cero \(\beta\in(z,a)\) de \(c_{2p}\). Como \(c_p\ne0\) en ese intervalo, el período crítico en \(\beta\) es exactamente \(2p\).

La analiticidad y \(c_{2p}(a)\ne0\) implican que los ceros de \(c_{2p}\) en \([\lambda_{n-1},a]\) son finitos. Entre sus ceros de período exacto \(2p\), situados a la derecha de \(\lambda_{n-1}\), existe así un primero. Es el selector original \(\lambda_n\), y satisface \(\lambda_n\le\beta<a\). \(\square\)

La utilización del último cero auxiliar \(z\) pertenece únicamente a la prueba de existencia. El parámetro seleccionado continúa siendo la \textbf{primera raíz nueva} \(\lambda_n\).


\subsection{Convergencia del selector global al primer itinerario infinito}
\label{hmtfg:clasificacion:5}

\textbf{Teorema de continuación global.} Para la familia HMT fijada y el selector original, todos los itinerarios seleccionados son canónicos y
\begin{equation}
\boxed{K_{\lambda_n}=W_n0,\qquad
\lambda_n\nearrow\min\Lambda_\tau.}
\label{hmtfg:clasificacion:eq:10}
\end{equation}

\textbf{Demostración.} Los lemas de comparación y existencia sitúan cada parámetro seleccionado por debajo de \(a\), con itinerario menor que \(\tau\). El teorema de clasificación identifica ese itinerario con \(W_n0\). La sucesión es creciente y acotada por \(a\); sea \(b\le a\) su límite.

Fijemos \(p\ge1\). Para \(n\) suficientemente grande, \(2p<2^n\). Entonces los signos de \(c_p(\lambda_n)\) y \(c_{2p}(\lambda_n)\) coinciden con \(\tau_p\) y \(\tau_{2p}=-\tau_p\), respectivamente. La cota uniforme de las segundas derivadas,
\[
\|(f_\lambda^p)''\|\le B_p,
\qquad B_p=16^p\frac{16^p-1}{15},
\]
y la identidad \((f_\lambda^p)'(0)=0\) implican
\[
|c_{2p}(\lambda_n)-c_p(\lambda_n)|
\le\tfrac12B_p|c_p(\lambda_n)|^2.
\]
Como los dos términos del lado izquierdo tienen signos contrarios,
\begin{equation}
|c_p(\lambda_n)|>2/B_p.
\label{hmtfg:clasificacion:eq:11}
\end{equation}
Por continuidad, \(c_p(b)\) conserva el signo \(\tau_p\) y satisface \(|c_p(b)|\ge2/B_p>0\). Esto vale para cada \(p\), de modo que \(b\in\Lambda_\tau\). La minimalidad de \(a\) da \(b\ge a\), y por tanto \(b=a\). \(\square\)

La prueba completa conserva el selector global y produce sus itinerarios y su límite. La información de los primeros ocho niveles es compatible con el resultado, pero la demostración de \eqref{hmtfg:clasificacion:eq:10} no extrapola esos ocho niveles: utiliza inducción, intervalos invariantes y separación uniforme de cada signo.


\subsection{Controles lógicos y falsadores}
\label{hmtfg:clasificacion:7}

\par\noindent\textbf{1.}\enspace El teorema de clasificación requiere un máximo unimodal estricto y un intervalo invariante. Si alguno de los dos ramos pierde monotonía, los intervalos \eqref{hmtfg:clasificacion:eq:2} pueden dejar de ser restrictivos y la clasificación debe comprobarse de nuevo.
\par\noindent\textbf{2.}\enspace La condición \(K<\tau\) es esencial. Después de la palabra límite pueden existir centros de período \(2^n\) con otras combinatorias; el teorema no los identifica con \(W_n0\).
\par\noindent\textbf{3.}\enspace La analiticidad en el parámetro proporciona aislamiento y finitud de los ceros sobre el compacto. Para una familia solamente continua, la existencia de una primera raíz estrictamente posterior requiere un argumento adicional.
\par\noindent\textbf{4.}\enspace Un límite de raíces críticas puede, en general, perder signos al alcanzar un retorno de período menor. La desigualdad \eqref{hmtfg:clasificacion:eq:11} excluye precisamente ese fenómeno para esta sucesión y cada horizonte fijo.
\par\noindent\textbf{5.}\enspace La convergencia de las primeras raíces al primer parámetro de itinerario \(\tau\) no equivale a la unicidad de todos los parámetros de ese itinerario ni determina por sí misma una tasa métrica. El cruce transversal y la identificación de la rama aportan esas conclusiones posteriores.
